\section{Conclusion}
\label{sec:conclusion}

IPI is a transport-independent application protocol for stateless CV/ITS
applications and stateful CAV applications. Its common interface allows both
workload classes to be evaluated over direct PC5 and network-assisted Uu rather
than assigning each class to one communication technology in advance. In a
tightly controlled real-vehicle deployment, the evaluation increases
application payload, reverses the direction of large CAV objects, combines
finite exchanges with sustained streams, and varies communication conditions.
It measures response latency, deadline completion, and exact directional
application goodput.

The results establish that the communication envelope is specific to each
application. Direct PC5 carries compact J2735 messages when roadside
coverage is available, while obstruction and sparse infrastructure reduce
their availability. The 5G path carries larger stateful exchanges, but the
measured uplink crosses a 500-ms decision deadline near 20~KiB. The GL-X3000
repeat measured approximately 31\% higher sustained goodput than the closest
MG52 control under the established profile. With the
GL-X3000 held constant, the deployed 40/40/20
configuration still produced more deadline misses and lower goodput than
70/20/10 despite its larger nominal uplink allocation. A continuing uplink
stream also moved compact foreground exchanges across their deadlines.

These results connect the two contributions. IPI expresses periodic CV
updates, stateful CAV bursts, and continuing CAV sessions through one
application interface. The communication study then identifies where payload,
direction, coverage, configuration, and simultaneous traffic prevent those
workloads from meeting their requirements. Supporting CAV applications
therefore requires radio vendors and mobile network operators to provide stable
bidirectional configurations and scheduling for both event-triggered bursts
and sustained deadline-sensitive streams.
